\documentclass[10pt,a4paper]{amsart}
\usepackage[margin=19mm]{geometry}
\usepackage{amsmath,amssymb,mathtools}
\usepackage[scr=rsfs,cal=euler]{mathalfa}
\usepackage{xfrac}
\usepackage[unicode,hidelinks,bookmarks=false]{hyperref}
\newcommand{\ie}{i.e.\ }
\newcommand{\cf}{cf.\ }
\newcommand{\resp}{respectively\ }
\newcommand{\Subsection}{\subsection}
\providecommand{\nobreakdash}{\nobreak\hbox{-}}
\setlength{\emergencystretch}{2em}
\multlinegap0pt
\usepackage{bm}
\usepackage{bbm}
\usepackage{dsfont}
\usepackage{xspace}
\usepackage{cancel}
\usepackage{mathtools}
\usepackage[shortlabels]{enumitem}
\usepackage{amsthm}
\makeatletter
\@ifundefined{proposition}{\newtheorem{proposition}{Proposition}}{}
\makeatother
\title[The one-point Schreier Poisson boundary of Thompson's group ]{The one-point Schreier Poisson boundary of Thompson's group $F$}
\author{Christian M\"onch}
\date{Source: arXiv:2606.23896v1 (CC BY 4.0)}
\begin{document}
\maketitle

\noindent\emph{Source and licence.} The one-point Schreier Poisson boundary of Thompson's group $F$. Version arXiv:2606.23896v1 (CC BY 4.0), distributed under the Creative Commons Attribution 4.0 International licence, as recorded in the arXiv licence metadata for this exact version and shown on the abstract page https://arxiv.org/abs/2606.23896v1. Full rights notice: licenses/arxiv-2606.23896-CC-BY-4.0.txt.\par\smallskip

\noindent\textit{Editorial scope.} Counted unit: the Proposition whose source label is \texttt{prop:trace-kernel} in the article (source0.tex lines 783--854, source label \texttt{prop:trace-kernel}), delivered together with the article's own complete original proof of it (source0.tex lines 855--996, one \texttt{proof} environment of 142 lines). No mathematics is written, repaired or re-derived: statement and proof are the article's own text line for line. Required context retained uncounted: none. Every definition, display and label of those lines is retained unchanged. Omitted from this package and not redistributed: 1--782, 997--1576. Nothing inside a retained excerpt is omitted or altered: every retained body is delivered exactly as the article's lines are written, so no omission receipt of any kind accompanies this package. Citations: the retained excerpts carry no citation command at all, so no bibliography entry of the article is transcribed and no reference is redistributed. Reference adaptation: none was needed and none was made; every reference inside the delivered unit resolves inside it, so no unresolved reference or citation is produced. Printed slips kept literal: no printed word, symbol, hypothesis or number of the counted statement or of its proof was altered. Layout and class: the article's own class is \texttt{scrartcl} with packages including iftex, mystyle-koma-enhanced, tikz; the wrapper is the lane's pinned \texttt{amsart} editorial wrapper at 10pt on A4, which restates the theorem-like environments the retained text needs and adds the article's own macro definitions that the retained text needs (none), with extra wrapper packages (bm, bbm, dsfont, xspace, cancel, mathtools, enumitem). The delivered numbering is the unit's own. Assets: no retained span contains a figure environment, a graphics-inclusion command, a PDF-page inclusion, a table or a TikZ picture; no image file is copied into this package, the package declares no asset dependency and no image file is redistributed with it. Licence: version arXiv:2606.23896v1 of this article is distributed under the Creative Commons Attribution 4.0 International licence, as recorded in the arXiv licence metadata for this exact version and shown on the abstract page https://arxiv.org/abs/2606.23896v1. A full rights notice is delivered at licenses/arxiv-2606.23896-CC-BY-4.0.txt. Characters: every retained excerpt is pure ASCII, so the delivered engine, XeLaTeX with its default Latin Modern text font, reports no missing character.
\begin{proposition}\label{prop:trace-kernel}
The no-holding trace \(Q\) is the reversible
nearest-neighbor random walk on the rooted binary tree \(B\) with conductances

\[
c(w,w0)=1,
\qquad
c(w,w1)=\frac12 .
\]

Equivalently, the edge resistances are

\[
r(w,w0)=1,
\qquad
r(w,w1)=2 .
\]

More explicitly, at the root,

\[
Q(\varnothing,0)=\frac23,
\qquad
Q(\varnothing,1)=\frac13 .
\]

If \(w\ne\varnothing\) ends in \(0\), then

\[
Q(w,w^-)=\frac25,
\qquad
Q(w,w0)=\frac25,
\qquad
Q(w,w1)=\frac15 .
\]

If \(w\ne\varnothing\) ends in \(1\), then

\[
Q(w,w^-)=\frac14,
\qquad
Q(w,w0)=\frac12,
\qquad
Q(w,w1)=\frac14 .
\]

The trace probabilities before deleting holds are:

\[
\begin{array}{c|c|c|c|c}
\text{state } \gamma
& \widehat P(\gamma,\gamma)
& \widehat P(\gamma,p(\gamma))
& \widehat P(\gamma,\gamma_0)
& \widehat P(\gamma,\gamma_1)
\\ \hline
\text{non-root }0\text{-child}
& 3/8 & 1/4 & 1/4 & 1/8
\\
\text{non-root }1\text{-child}
& 1/2 & 1/8 & 1/4 & 1/8
\\
o=\varnothing
& 5/8 & - & 1/4 & 1/8
\end{array}
\]

Here \(p(\gamma)\) denotes the parent of a non-root vertex, and
\(\gamma_i=f_i(\gamma)\). The no-holding kernel \(Q\) is obtained by
normalizing each row after removing the holding column.

\end{proposition}

\begin{proof}
Compute \(\widehat P\) before deleting holds. Let
\(\gamma\in B\), and set

\[
\gamma_0=f_0(\gamma),
\qquad
\gamma_1=f_1(\gamma).
\]

Every non-root grey vertex lies in exactly one of the two intervals
\((0,1/2)\) and \((1/2,1)\). These are precisely the \(0\)-children and
\(1\)-children, respectively. Indeed, if \(0<\gamma<1/2\), then its parent is
\(2\gamma\), while if \(1/2<\gamma<1\), then its parent is \(2\gamma-1\).

First assume \(0<\gamma<1/2\). Then \(\gamma\) is a \(0\)-child of its parent,
so

\[
p(\gamma)=2\gamma.
\]

The four generator moves behave as follows:

\begin{itemize}[leftmargin=*]
\item \(\widetilde A(\gamma)=\gamma-1\) enters a negative recurrent ray and
  returns to \(\gamma\) before the next hit of \(B\).
\item \(\widetilde A^{-1}(\gamma)=\gamma+1\) enters the white vertex between
  \(\gamma\) and \(\gamma_1\), from which the trace hits these two grey
  vertices with probabilities \(1/2\) and \(1/2\).
\item \(\widetilde B(\gamma)=\gamma_0\).
\item \(\widetilde B^{-1}(\gamma)=2\gamma=p(\gamma)\).
\end{itemize}

Therefore

\[
\widehat P(\gamma,\gamma)=\frac14+\frac14\cdot\frac12=\frac38,
\]

\[
\widehat P(\gamma,\gamma_0)=\frac14,
\qquad
\widehat P(\gamma,\gamma_1)=\frac14\cdot\frac12=\frac18,
\qquad
\widehat P(\gamma,p(\gamma))=\frac14 .
\]

Deleting the hold gives

\[
Q(\gamma,p(\gamma))=\frac{1/4}{5/8}=\frac25,
\qquad
Q(\gamma,\gamma_0)=\frac{1/4}{5/8}=\frac25,
\qquad
Q(\gamma,\gamma_1)=\frac{1/8}{5/8}=\frac15 .
\]

Now assume \(1/2<\gamma<1\). Then \(\gamma\) is a \(1\)-child and its parent is

\[
p(\gamma)=2\gamma-1.
\]

The move \(\widetilde B^{-1}(\gamma)=2\gamma\) enters the white vertex between
\(p(\gamma)\) and \(\gamma\), so it contributes half to holding and half to
the parent. The same local computation gives

\[
\widehat P(\gamma,\gamma)=\frac12,
\qquad
\widehat P(\gamma,\gamma_0)=\frac14,
\qquad
\widehat P(\gamma,\gamma_1)=\frac18,
\qquad
\widehat P(\gamma,p(\gamma))=\frac18 .
\]

Deleting holds gives

\[
Q(\gamma,p(\gamma))=\frac14,
\qquad
Q(\gamma,\gamma_0)=\frac12,
\qquad
Q(\gamma,\gamma_1)=\frac14 .
\]

At the root \(o=1/2\), the four generator moves are:

\begin{itemize}[leftmargin=*]
\item \(\widetilde A(o)=-1/2\), which enters the negative type ray attached to
  \(o\) and returns to \(o\) almost surely before the next hit of \(B\).
\item \(\widetilde A^{-1}(o)=3/2\), a white vertex whose first hit of \(I\) is
  \(o\) or \(3/4=o1\), each with probability \(1/2\).
\item \(\widetilde B(o)=1/4=o0\).
\item \(\widetilde B^{-1}(o)=1\). From \(1\), the walk hits \(o\) almost surely
  before any other point of \(B\): excursions from \(1\) to \(0\) or into the
  positive type ray are recurrent, and the edge \(1\to o\) is available at
  each return to \(1\).
\end{itemize}

Therefore

\[
\widehat P(o,o)=\frac58,
\qquad
\widehat P(o,o0)=\frac14,
\qquad
\widehat P(o,o1)=\frac18,
\]

and therefore

\[
Q(o,o0)=\frac{1/4}{3/8}=\frac23,
\qquad
Q(o,o1)=\frac{1/8}{3/8}=\frac13 .
\]

These probabilities are those of the stated conductance walk. If
\(w\) ends in \(0\), then the parent edge has conductance \(1\), while the two
child edges have conductances \(1\) and \(1/2\). Normalizing gives

\[
\frac{1}{1+1+1/2}=\frac25,
\qquad
\frac{1/2}{1+1+1/2}=\frac15 .
\]

If \(w\) ends in \(1\), then the parent edge has conductance \(1/2\), while
the two child edges have conductances \(1\) and \(1/2\). Normalizing gives

\[
\frac{1/2}{1/2+1+1/2}=\frac14,
\qquad
\frac{1}{1/2+1+1/2}=\frac12 .
\]

The root has incident conductances \(1\) and \(1/2\), giving probabilities
\(2/3\) and \(1/3\).
\end{proof}

\end{document}
